\section{Antecedents and neighbouring programmes: digital physics and computational views of nature}
\label{sec:digital-physics}

The simulation hypothesis is often presented as the latest form of an older
idea, that ``the universe is a computer''. That older literature contains
several different theses. If they are run together, it becomes unclear what a
given observation could test. The key distinction for this review is between
the claim that the universe \emph{is} a computation (or is computable, or is
discrete) and the claim that it is \emph{run on} a computer that belongs to
some other reality. This section traces where the computational views came
from and fixes the terms used in later sections.

\subsection{Six theses that are often conflated}
\label{sec:digital-physics:theses}

We separate six claims. Table~\ref{tab:digital-physics:theses} summarizes them.

\emph{(C) Computational modellability.} Floridi separates two questions. The
first is whether the physical universe can be adequately modelled digitally
and computationally. The second is whether it is digital and computational in
itself. He treats the first as an open ``empirico-mathematical'' question and
argues that the second is ill-posed~\cite{Floridi2009against}. Thesis (C)
carries no ontological commitment. Floridi puts this as a contrast between the
``predicative'' and ``attributive'' senses of ``digital physics''.

\emph{(CU) Computable universe.} Szudzik defines a computable physical model as
a recursive set of states together with a total recursive function for each
observable. On this basis he states a \emph{computable universe hypothesis}:
the universe has a recursive set of states, and every observable is a total
recursive function of the state~\cite{Szudzik2012computable}. His examples
include a model with non-discrete, continuous motion. So computability in this
sense does not require discreteness. Tegmark's \emph{Computable Universe
Hypothesis} is a different formulation. It requires that the relations defining
the mathematical structure identified with our universe be computable by
halting computations~\cite{Tegmark2008mathematical}. Deutsch's physical
reading of the Church--Turing thesis is a related but separate claim. It says
that ``every finitely realizable physical system can be perfectly simulated by
a universal model computing machine operating by finite
means''~\cite{Deutsch1985quantum}. Deutsch argued that classical continuum
physics does not satisfy this strong form with respect to the Turing machine.
Quantum theory, he argued, is compatible with it through a universal quantum
computer. He also held that the principle has the epistemological status of
other physical principles, such as the third law of
thermodynamics~\cite{Deutsch1985quantum}. Deutsch himself calls it ``the
Church--Turing principle''.

\emph{(D) Digital universe.} This is the thesis that the physical universe
(time, space and every entity and process in spacetime) is ultimately discrete.
Floridi takes it as the core of ``digital ontology''~\cite{Floridi2009against}.
The SEP entry on physical computation adds a requirement for a
cellular-automaton universe. All fundamental magnitudes must take finitely many
values, and space and time must be fundamentally discrete or emerge from
discrete processing~\cite{Piccinini2025computation}.

\emph{(P) Ontic pancomputationalism.} Here the universe is itself a computing
system. The SEP entry notes that this view originated in physics, not
philosophy. It combines an \emph{empirical} claim with a \emph{metaphysical}
one. The empirical claim is that fundamental magnitudes and their state
transitions are \emph{exactly} described by some computational formalism. The
metaphysical claim is that computation is what the universe is made
of~\cite{Piccinini2025computation}. There is a classical version, based on
cellular automata, and a quantum version associated with
Lloyd~\cite{Lloyd2006programming}. The entry describes the quantum version as
possibly a reformulation of quantum mechanics without change in empirical
content~\cite{Piccinini2025computation}.

\emph{(M) Mathematical universe.} Tegmark's Mathematical Universe Hypothesis
states that ``our physical world is an abstract mathematical
structure''~\cite{Tegmark2008mathematical}. It generalizes his earlier
``ultimate ensemble'' postulate that all structures that exist mathematically
also exist physically~\cite{Tegmark1998is}.

\emph{(S) Simulated universe.} This is the claim that our universe is a
computation \emph{implemented on hardware that exists outside it}. Usually the
hardware is also assumed to have been built by designers. This is the
simulation hypothesis proper, whose philosophical formulation is reviewed in
the section on the simulation argument (Sec.~\ref{sec:philosophy}).


\begin{table}[t]
\centering
\small
\caption{Computational theses about the physical world. The last column states
whether the thesis itself posits computing hardware outside our universe.}
\label{tab:digital-physics:theses}
\begin{tabular}{@{}p{0.19\linewidth}p{0.36\linewidth}p{0.25\linewidth}p{0.10\linewidth}@{}}
\hline
Thesis & Core claim & Representative sources & External hardware? \\
\hline
(C) Modellability & The world can be adequately modelled computationally & \cite{Floridi2009against} & No \\
(CU) Computable universe & States and observables (or defining relations) are computable & \cite{Szudzik2012computable,Tegmark2008mathematical,Deutsch1985quantum} & No \\
(D) Digital universe & Spacetime and all magnitudes are fundamentally discrete & \cite{Floridi2009against,Fredkin2003introduction} & No \\
(P) Ontic pancomputationalism & The universe is a (classical or quantum) computing system & \cite{Piccinini2025computation,Lloyd2002computational,Wolfram2002new} & Not necessarily \\
(M) Mathematical universe & The world is a mathematical structure & \cite{Tegmark2008mathematical} & No \\
(S) Simulated universe & The world is a computation run on hardware outside it & \cite{Moravec1992pigs,Bostrom2003are} & Yes \\
\hline
\end{tabular}
\end{table}

Several relations among these theses are stated explicitly in the literature.
First, (D) and (P) are independent. Floridi cites Wheeler as someone who held
the former but not, at least not explicitly, the latter~\cite{Floridi2009against}.
Second, (CU) does not entail (D), as Szudzik's continuous example
shows~\cite{Szudzik2012computable}. Third, (P) does not entail (S). The SEP
entry notes that the metaphysical claim of (P) needs an account of what the
computation \emph{is}. It lists three options: a simulation running on a
computer in another universe, computations as abstract mathematical entities
(a computational Pythagoreanism), and computational structural
realism~\cite{Piccinini2025computation}. Only the first is (S). Fourth,
Tegmark contrasts his hypotheses with (S) explicitly. The Computable Universe
Hypothesis ``requires the description (the relations) rather than the time
evolution to be computable'', so that ``the universe is a mathematical
structure according to the MUH, as opposed to a computation according to the
simulation hypothesis''~\cite{Tegmark2008mathematical}.

\subsection{A genealogy}
\label{sec:digital-physics:genealogy}

\paragraph{Zuse and Fredkin.} Zuse's ``Rechnender Raum'' first appeared as a
journal article in 1967~\cite{Zuse1967rechnender}. It was followed by a book in
1969~\cite{Zuse1969rechnender} and an MIT technical translation, \emph{Calculating
Space}, in 1970~\cite{Zuse1970calculating}. We could not consult Zuse's text,
so we rely on secondary characterizations. The SEP entry credits Zuse and
Edward Fredkin with the ``earliest and best known'' version of ontic
pancomputationalism, in which the universe is a giant cellular automaton. It
adds that Fredkin's largely unpublished ideas influenced Feynman, Toffoli and
Wolfram~\cite{Piccinini2025computation}. Floridi calls the shared position the
``Zuse Thesis'': the universe is deterministically computed on a giant but
discrete computer. He notes that Zuse, Fredkin and Wolfram took that computer
to be a cellular automaton, Schmidhuber a universal Turing machine, and Lloyd a
quantum computer~\cite{Floridi2009against}. Fredkin developed ``digital
mechanics''~\cite{Fredkin1990informational} and ``digital
philosophy''~\cite{Fredkin2003introduction}. As quoted by Floridi, digital
philosophy takes space and time to be discrete and the bit as the unit of
state. It also demands the ``eventual ability to derive'' the equations of
current physics from such models~\cite{Floridi2009against}. According to the
SEP entry, Fredkin~\cite{Fredkin2003introduction} also considered the
possibility that our universe is a simulation on a computer in another universe
whose physics is inaccessible to us~\cite{Piccinini2025computation}. That
possibility is thesis (S), not (D).

\paragraph{The MIT physics-of-computation papers (1981--1984).} Feynman's
keynote ``Simulating physics with computers'' says that his interest in the
topic was inspired by Fredkin~\cite{Feynman1982simulating}. The part that
concerns us here is not complexity, which is treated in the section on
computational complexity (Sec.~\ref{sec:complexity}),
{} but Feynman's notion of \emph{exact} simulation. Exact simulation would
require ``everything that happens in a finite volume of space and time'' to be
``exactly analyzable with a finite number of logical operations''. Present
physics does not meet this condition, so ``if this proposition is right,
physical law is wrong''~\cite{Feynman1982simulating}. Feynman also noted that
if space were replaced by a simple lattice, the speed of light would generically
depend slightly on direction. The resulting anisotropies are experimentally
constrained~\cite{Feynman1982simulating}. This observation anticipates the
lattice tests reviewed in the section on lattice signatures and Lorentz
violation (Sec.~\ref{sec:lattice-liv}).
He called Fredkin's programme of finding a computer simulation of physics ``an
excellent program'', provided the simulation is quantum
mechanical~\cite{Feynman1982simulating}. Fredkin and Toffoli's conservative
logic went the other way, from computation to physics. It is a model of
computation built to respect reversibility and conservation laws, and it shows
that a perfect gas in a suitably shaped container can reproduce a
general-purpose digital computer~\cite{Fredkin1982conservative}. Toffoli later
proposed cellular automata ``as an alternative to (rather than an approximation
of)'' differential equations in modelling physics~\cite{Toffoli1984cellular}.
Zuse~\cite{Zuse1982computing} and Wheeler~\cite{Wheeler1982computer} also
published papers in the 1982 issue of the \emph{International Journal of
Theoretical Physics} that carried Feynman's keynote.

\paragraph{Wheeler: ``it from bit''.} Wheeler's slogan is often grouped with
digital physics, but its content is different. It holds that every physical
quantity ``derives its ultimate significance from bits, binary yes-or-no
indications''. These bits are elementary acts of
``observer-participancy''~\cite{Wheeler1990information}. Wheeler denied a
microscopic continuum of space and time (``continuum-based physics, no;
information-based physics, yes''). He also explicitly rejected ``the concept of
universe as machine'', partly because it requires postulating ``a
supermachine''~\cite{Wheeler1990information}. ``It from bit'' is therefore an
informational and participatory thesis close to (D). It is not a thesis of
type (P) or (S).

\paragraph{Lloyd: limits and the computing universe.} Lloyd derived that a
system of mean energy $E$ can perform at most $2E/\pi\hbar$ elementary
operations per second. For a 1\,kg ``ultimate laptop'' the bound is about
$5.4\times10^{50}$ operations per second~\cite{Lloyd2000ultimate}. Applied to
the observable universe, the same limits give order-of-magnitude upper bounds
of about $10^{120}$ operations on about $10^{90}$ bits (about $10^{120}$ bits if
gravitational degrees of freedom are included). The operation count scales as
$(t/t_P)^2$~\cite{Lloyd2002computational}. Lloyd offered three readings of
these numbers. They bound the computation that matter can have performed. They
bound from below the resources a quantum computer needs to simulate the
universe directly. And, in a reading he calls ``more controversial'', they give
the memory and operation count of the universe regarded as a
computation~\cite{Lloyd2002computational}. The second reading bears on (S).
The third is a quantum version of (P). Lloyd's own answer to ``Is the universe
a computer?'' depends ``on the meaning of `computer' and on the meaning of
`is'\,'': the universe is not a conventional digital computer, but ``it can
still compute''~\cite{Lloyd2002computational}. A later preprint proposes
deriving spacetime geometry from an underlying quantum
computation~\cite{Lloyd2005theory}.

\paragraph{Wolfram: \emph{A New Kind of Science} and the 2020 models.} Wolfram
asked whether ``some simple program which, if run for long enough, would
reproduce our universe in every detail'' might underlie physics. Such a program
would be a model ``not in any sense an approximation or
idealization''~\cite{Wolfram2002new}. In the same book he wrote that no literal
mechanistic model of the Zuse--Fredkin kind ``can ever in the end realistically
be expected to work''~\cite[p.~1026]{Wolfram2002new}. Expert reviews of the
physics chapter were critical. Aaronson argued that its causal-network picture
of spacetime had been discussed earlier in quantum gravity, with the difference
that Wolfram's model is explicitly classical. He also showed that Wolfram's
deterministic ``long-range thread'' account of entanglement cannot satisfy
causal invariance, the relativity postulate and Bell-inequality violation all
at once~\cite{Aaronson2002book}. Kadanoff attributed the automaton-universe view
to Fredkin. He judged the chapter ``a partially formed idea---exciting, but not
yet science'' and found ``no comparison with experiment''~\cite{Kadanoff2002new}.
Weinberg's essay review saw no motivation for the discrete-universe speculation
beyond familiarity with computers. He also noted that the computational
equivalence of automata to continuous systems holds only if nothing in nature is
truly continuous~\cite{Weinberg2002is}. In 2020 Wolfram introduced models in
which the universe is an evolving hypergraph with no intrinsic space, time or
matter. He proposed that ``causal invariance'' of the rewriting rules yields
Lorentz invariance, general covariance and local gauge invariance, and
described parts of this correspondence as ``speculative''~\cite{Wolfram2020class}.
In that paper experimentally accessible predictions are a prospect (``it seems
increasingly likely that experimentally accessible predictions will be
possible''), not a result. The models' elementary length need not coincide with
the Planck length~\cite{Wolfram2020class}. Gorard reports proofs that causal
invariance is equivalent to a discrete general covariance, and derives discrete
analogues of Lorentz covariance and, under a dimension-preservation assumption,
of the Einstein equations~\cite{Gorard2020some}. He also relates the models to
causal sets~\cite{Gorard2020algorithmic}. Butterfield and Dowker argue that any
Planck-scale discrete theory that recovers general relativity must contain
causal sets. They list the Wolfram model among the approaches whose status
under this argument ``it would be a good project to
assess''~\cite{Butterfield2024recovering}. The 2020 models allow multiway
(branching) evolution, which the NKS framework did not. We did not find a
published analysis of whether Aaronson's argument carries over to them.

\paragraph{'t Hooft: the cellular automaton interpretation.} 't~Hooft treats
quantum mechanics ``as a tool, not as a theory''. He suggests that the Standard
Model with gravity may be a quantum-mechanical analysis of a system that
``could be classical at its core''. He also argues that the usual objections to
superdeterminism can be overcome ``at least in principle''~\cite{tHooft2016cellular}.
This is a (D)-type proposal about the foundations of quantum mechanics. It does
not posit a simulator. It makes a ``firm prediction'': quantum computers will
not outperform a classical computer with one memory site per Planck volume (or
Planck area) that performs one operation per Planck time~\cite{tHooft2016cellular}.
This links the proposal to the complexity questions treated in the section on
computational complexity (Sec.~\ref{sec:complexity}).


\paragraph{Algorithmic ensembles and the mathematical universe.} Schmidhuber
cast his 1997 analysis explicitly as a simulation: a ``Great Programmer''
runs all computable universes on ``His Big Computer''. Here ``computable''
means discrete time and finitely describable states, and computing all
universes needs far less information than computing one arbitrary
universe~\cite{Schmidhuber1997computer}. He noted that inhabitants ``run on
local time'' and cannot tell how many machine steps each of their time steps
takes~\cite{Schmidhuber1997computer}. His later ``speed prior'' assigns low
probability to histories that are slow to compute. From it he derived a
concrete prediction: apparently random events such as
individual beta decays come from a fast pseudorandom
generator~\cite{Schmidhuber2000algorithmic}. Tegmark separates his
computable-universe hypothesis from (S), as quoted above. He also calls
``a common misconception in the universe simulation literature'' the
identification of physical time with the step-by-step flow of a computation.
On his account, a simulating computer need only \emph{specify} a history, not
compute it~\cite{Tegmark2008mathematical}. He concedes that ``virtually all
historically successful theories of physics violate the CUH'' because they use
the continuum~\cite{Tegmark2008mathematical}. Deutsch's and Feynman's
observations point to the same tension~\cite{Deutsch1985quantum,Feynman1982simulating}.

\paragraph{Early simulated-world proposals.} In a 1992 essay Moravec argued
that a future cyberspace would replay human history many times. Hence ``the very moment we are now
experiencing may actually be (almost certainly is)'' such a replayed
event~\cite{Moravec1992pigs}. This informally anticipates the counting
reasoning of the simulation argument~\cite{Bostrom2003are}. Bostrom's paper cites Moravec's
\emph{Mind Children}~\cite{Moravec1988mind} and Lloyd's limits
paper~\cite{Lloyd2000ultimate}. Tegmark lists Tipler~\cite{Tipler1994physics},
Bostrom and Schmidhuber as authors who discussed the probability that we are
simulated~\cite{Tegmark2008mathematical}. He describes Lloyd's view as an
intermediate case: an analog quantum simulation ``not designed by
anybody''~\cite{Tegmark2008mathematical}.

\subsection{Philosophical assessments}
\label{sec:digital-physics:philosophy}

The SEP entry distinguishes ``unlimited'' pancomputationalism from ``limited''
pancomputationalism. The unlimited form, associated with Putnam and Searle,
holds that every sufficiently complex system implements very many computations.
The limited form holds that every system implements some computation. Both
differ from the ontic pancomputationalism of the physicists discussed
here~\cite{Piccinini2025computation}. The entry finds ``little positive
evidence'' for ontic pancomputationalism. It suggests that supporters are
motivated by ``the desire for exact computational models of the world'' rather
than by evidence. It notes the metaphysical objection that computational
entities lack the causal and qualitative properties we observe. It also grants
that the classical version is ``in principle'' testable~\cite{Piccinini2025computation}.
Floridi argues that digital ontology should be abandoned in favour of
informational structural realism. His reason is that ``digital'' and
``analogue'' are modes of presentation relative to a level of abstraction, not
properties of reality in itself~\cite{Floridi2009against}. His argument has
drawn at least one published reply~\cite{Sdrolia2014rethinking}.

\subsection{Consequences for this review}
\label{sec:digital-physics:consequences}

These distinctions have three practical consequences for the rest of the
review. The first two are our own inferences from the sources above; the third
restates positions in them.
First, empirical searches for spacetime discreteness (section on lattice
signatures) test particular (D)-type models, such as a lattice with preferred directions. Their
bearing on (S) depends on extra assumptions about how a simulator would
represent our world. Nothing in (S) as such fixes a lattice representation.
It has also been argued that the simulator's steps need not correspond to our
time, and that a simulator might only specify our history rather than compute
it step by step~\cite{Schmidhuber1997computer,Tegmark2008mathematical}.
Second, the computability and complexity results (section on complexity) bear
on (CU) and on the resource side of (S). Examples are Deutsch's principle, Lloyd's bounds
(read as lower bounds on the cost of a direct quantum simulation) and
't~Hooft's prediction about quantum
computers~\cite{Deutsch1985quantum,Lloyd2002computational,tHooft2016cellular}.
Third, arguments that the universe ``is'' computational do not by themselves
support the claim that it is simulated, because a computational ontology can be
read in ways that posit no external hardware~\cite{Piccinini2025computation,Tegmark2008mathematical}.
